% ============================================================
%  يُترجَم بـ XeLaTeX  (في Overleaf: Menu → Compiler → XeLaTeX)
%  الخط: Amiri (متوفر في Overleaf، أو حمّله من github.com/aliftype/amiri)
% ============================================================
\documentclass[12pt,a4paper]{article}
\usepackage[a4paper,margin=1.7cm,top=2.3cm,bottom=2.9cm,headheight=15pt,footskip=1.5cm]{geometry}
\usepackage{amsmath,amssymb}
\usepackage{icomma}
\usepackage[english,bidi=default,layout=lists]{babel}
\babelprovide[import,main]{arabic}
\babelfont{rm}[Scale=1.08]{Amiri}
\babelfont{sf}[Scale=1.08]{Amiri}
\usepackage[most]{tcolorbox}
\usepackage{tikz}
\usetikzlibrary{arrows.meta,calc,shadings}
\usepackage{array}
\usepackage{colortbl}
\usepackage{fancyhdr}
\usepackage{titlesec}
\usepackage{enumitem}
\usepackage{collcell}
\usepackage{needspace}
\usepackage{etoolbox}
\newcolumntype{C}{>{\collectcell\ar}c<{\endcollectcell}}
\pretocmd{\section}{\Needspace{15\baselineskip}}{}{}
\pretocmd{\subsection}{\Needspace{11\baselineskip}}{}{}

\definecolor{cDef}{HTML}{1F4E9C}
\definecolor{cThm}{HTML}{B3261E}
\definecolor{cRem}{HTML}{2E7D32}
\definecolor{cEx}{HTML}{C76B00}
\definecolor{cMeth}{HTML}{6A1B9A}
\definecolor{cSol}{HTML}{00695C}
\definecolor{cSec}{HTML}{1F3A68}

\titleformat{\section}{\Large\bfseries\color{cSec}}{\thesection.}{0.6em}{}[{\color{cSec!60}\titlerule[1pt]}]
\titleformat{\subsection}{\large\bfseries\color{cThm!85!black}}{\thesubsection.}{0.5em}{}
\titlespacing*{\section}{0pt}{14pt}{8pt}
\titlespacing*{\subsection}{0pt}{10pt}{5pt}
\setlength{\parindent}{0pt}
\setlength{\parskip}{4pt}
\setlist{itemsep=2pt,topsep=3pt}

\newtcolorbox{arbox}[3][]{enhanced,
  colback=#2!5,colframe=#2!80!black,arc=2mm,boxrule=0.8pt,
  left=6pt,right=6pt,top=8pt,bottom=5pt,
  fonttitle=\bfseries,coltitle=white,
  title={\foreignlanguage{arabic}{#3}},
  attach boxed title to top right={xshift=-5mm,yshift=-3mm},
  boxed title style={colback=#2!80!black,arc=1.5mm,boxrule=0pt},
  before upper={\selectlanguage{arabic}},
  before skip=8pt,after skip=8pt,breakable,#1}
\newenvironment{BOX}[2]{\begin{otherlanguage}{english}\begin{arbox}{#1}{#2}}{\end{arbox}\end{otherlanguage}}
\newenvironment{tarif}[1][تعريف]{\begin{BOX}{cDef}{#1}}{\end{BOX}}
\newenvironment{khassia}[1][خاصية]{\begin{BOX}{cThm}{#1}}{\end{BOX}}
\newenvironment{molahada}[1][ملاحظة]{\begin{BOX}{cRem}{#1}}{\end{BOX}}
\newenvironment{mithal}[1][مثال]{\begin{BOX}{cEx}{#1}}{\end{BOX}}
\newenvironment{tariqa}[1][طريقة]{\begin{BOX}{cMeth}{#1}}{\end{BOX}}

\newcounter{exo}
\newenvironment{tamrin}[1]{\stepcounter{exo}\begin{BOX}{cSec}{التمرين \theexo: #1}}{\end{BOX}}
\newenvironment{hal}{\Needspace{12\baselineskip}\begin{otherlanguage}{english}\begin{arbox}[breakable]{cSol}{حل التمرين \theexo}}{\end{arbox}\end{otherlanguage}}

\newenvironment{LTR}{\begin{otherlanguage}{english}}{\end{otherlanguage}}
\newcommand{\ar}[1]{\foreignlanguage{arabic}{#1}}

\renewcommand{\dots}{\mathinner{\ldotp\ldotp\ldotp}}
\renewcommand{\ldots}{\mathinner{\ldotp\ldotp\ldotp}}
\renewcommand{\cdots}{\mathinner{\cdotp\cdotp\cdotp}}
\newcommand{\N}{\mathbb{N}}
\newcommand{\Z}{\mathbb{Z}}
\newcommand{\D}{\mathbb{D}}
\newcommand{\Q}{\mathbb{Q}}
\newcommand{\R}{\mathbb{R}}
\newcommand{\pgcd}{\operatorname{PGCD}}
\newcommand{\ppcm}{\operatorname{PPCM}}

% \NLseg{a}{b}{[}{]} ; \NLright{a}{[} ; \NLleft{b}{]}  (interval on a number line)
\newcommand{\NLseg}[4]{\begin{LTR}\begin{tikzpicture}[baseline=-0.6ex]
 \draw[-{Stealth[length=2mm]},gray!70](0,0)--(3.4,0);
 \draw[line width=2pt,cThm!80!black](0.7,0)--(2.7,0);
 \node[cThm!80!black,font=\LARGE] at (0.7,0){$#3$};
 \node[cThm!80!black,font=\LARGE] at (2.7,0){$#4$};
 \node[below=3pt,font=\footnotesize] at (0.7,0){$#1$};
 \node[below=3pt,font=\footnotesize] at (2.7,0){$#2$};
\end{tikzpicture}\end{LTR}}
\newcommand{\NLright}[2]{\begin{LTR}\begin{tikzpicture}[baseline=-0.6ex]
 \draw[-{Stealth[length=2mm]},gray!70](0,0)--(3.4,0);
 \draw[line width=2pt,cThm!80!black](1.3,0)--(3.3,0);
 \node[cThm!80!black,font=\LARGE] at (1.3,0){$#2$};
 \node[below=3pt,font=\footnotesize] at (1.3,0){$#1$};
\end{tikzpicture}\end{LTR}}
\newcommand{\NLleft}[2]{\begin{LTR}\begin{tikzpicture}[baseline=-0.6ex]
 \draw[-{Stealth[length=2mm]},gray!70](0,0)--(3.4,0);
 \draw[line width=2pt,cThm!80!black](0.05,0)--(2.1,0);
 \node[cThm!80!black,font=\LARGE] at (2.1,0){$#2$};
 \node[below=3pt,font=\footnotesize] at (2.1,0){$#1$};
\end{tikzpicture}\end{LTR}}

% ---------- social-media icons (pure TikZ) ----------
\definecolor{cFB}{HTML}{1877F2}
\definecolor{cTG}{HTML}{229ED9}
\definecolor{cIG1}{HTML}{FEDA75}
\definecolor{cIG2}{HTML}{D62976}
\definecolor{cIG3}{HTML}{4F5BD5}
\newcommand{\fbicon}{\begin{LTR}\begin{tikzpicture}[baseline=-0.75ex,x=1em,y=1em]
 \fill[cFB] (0,0) circle (0.55);
 \node[white,inner sep=0,font=\usefont{TU}{lmss}{bx}{n}\small] at (0.04,-0.02) {f};
\end{tikzpicture}\end{LTR}}
\newcommand{\igicon}{\begin{LTR}\begin{tikzpicture}[baseline=-0.75ex,x=1em,y=1em]
 \shade[lower left=cIG1,lower right=cIG2,upper left=cIG2,upper right=cIG3,rounded corners=0.28em] (-0.55,-0.55) rectangle (0.55,0.55);
 \draw[white,line width=0.09em,rounded corners=0.18em] (-0.36,-0.36) rectangle (0.36,0.36);
 \draw[white,line width=0.09em] (0,0) circle (0.17);
 \fill[white] (0.22,0.22) circle (0.05);
\end{tikzpicture}\end{LTR}}
\newcommand{\tgicon}{\begin{LTR}\begin{tikzpicture}[baseline=-0.75ex,x=1em,y=1em]
 \fill[cTG] (0,0) circle (0.55);
 \fill[white] (-0.33,0.0) -- (0.32,0.27) -- (0.2,-0.3) -- (0.0,-0.15) -- (-0.08,-0.28) -- (-0.08,-0.08) -- cycle;
 \fill[cTG!70!black] (-0.08,-0.08) -- (0.25,0.2) -- (0.0,-0.15) -- (-0.08,-0.28) -- cycle;
\end{tikzpicture}\end{LTR}}
% page number: black box, white two-digit number, at the right of the social row
\newcommand{\PageBox}{\begin{LTR}\tikz[baseline=(p.base)]\node[fill=black,text=white,inner sep=2.5pt,minimum width=1.6em,minimum height=1.3em,font=\usefont{TU}{lmss}{bx}{n}\small](p){\ifnum\value{page}<10 0\fi\arabic{page}};\end{LTR}}
% visual left-to-right order: Telegram | Instagram | Facebook | page box
\newcommand{\SocialRow}{\begin{otherlanguage}{english}%
 \ar{\textbf{الأستاذ سوايسية هشام للرياضيات}}\ \tgicon\hfill
 \textbf{prof\_hicham\_math}\ \igicon\hfill
 \ar{\textbf{الأستاذ سوايسية هشام للرياضيات}}\ \fbicon\hspace{2.2em}\PageBox\end{otherlanguage}}

\pagestyle{fancy}
\fancyhf{}
\renewcommand{\headrulewidth}{0.6pt}
\renewcommand{\footrulewidth}{0.4pt}
% \SetHeader{right head: document title}{left head: level -- stream}
\newcommand{\SetHeader}[2]{%
 \fancyhead[C]{\parbox{\textwidth}{\ar{#1\hfill #2}}}%
 \fancyfoot[C]{\parbox{\textwidth}{\small\SocialRow}}}
\usepackage{tkz-tab}
\usepackage{multicol}
% variation / sign table (LTR)
\newenvironment{TAB}{\par\begin{center}\begin{LTR}\begin{tikzpicture}[scale=0.95,every node/.style={font=\small}]}{\end{tikzpicture}\end{LTR}\end{center}}
% grid + axes, clipped: \begin{GR}{xmin}{xmax}{ymin}{ymax}{unit cm}
\newenvironment{GR}[5]{\par\begin{center}\begin{LTR}\begin{tikzpicture}[x=#5cm,y=#5cm,font=\footnotesize]
 \draw[gray!25,very thin] (#1,#3) grid (#2,#4);
 \draw[-{Stealth[length=2mm]}] (#1,0)--(#2+0.3,0) node[below left]{$x$};
 \draw[-{Stealth[length=2mm]}] (0,#3)--(0,#4+0.3) node[below left]{$y$};
 \node[below left] at (0,0){$O$};
 \begin{scope}\clip (#1,#3) rectangle (#2,#4);}
{\end{scope}\end{tikzpicture}\end{LTR}\end{center}}
\newcommand{\Ticks}[2]{\foreach \t in {#1} \node[below,font=\scriptsize] at (\t,0){$\t$}; \foreach \t in {#2} \node[left,font=\scriptsize] at (0,\t){$\t$};}
\newcommand{\bx}[1]{\ensuremath{\colorbox{cDef!12}{$\displaystyle #1$}}}
\newcommand{\TwoCol}[4]{\par\begin{LTR}\noindent\begin{minipage}[t]{#1\textwidth}\selectlanguage{arabic}#2\end{minipage}\hfill\begin{minipage}[t]{#3\textwidth}\selectlanguage{arabic}#4\end{minipage}\end{LTR}\par}
\newcommand{\RL}[4]{\TwoCol{#3}{#4}{#1}{#2}}  % \RL{w1}{right content}{w2}{left content}
\newcommand{\ok}{$\textcolor{cRem}{\checkmark}$}
\newcommand{\limn}{\displaystyle\lim_{n\to+\infty}}
% sub-questions, each on its own line
\newenvironment{sq}{\begin{itemize}[leftmargin=2.4em,itemsep=3pt]}{\end{itemize}}
\newcommand{\qa}{\item[(أ)]}\newcommand{\qb}{\item[(ب)]}\newcommand{\qc}{\item[(جـ)]}\newcommand{\qd}{\item[(د)]}\newcommand{\qe}{\item[(هـ)]}
\newcommand{\pa}{\par\smallskip\textbf{(أ)} }\newcommand{\pb}{\par\smallskip\textbf{(ب)} }\newcommand{\pc}{\par\smallskip\textbf{(جـ)} }\newcommand{\pd}{\par\smallskip\textbf{(د)} }\newcommand{\pe}{\par\smallskip\textbf{(هـ)} }
\newcommand{\QQ}[1]{\par\smallskip\textbf{#1.} }
\newenvironment{tmr}{\stepcounter{exo}\begin{BOX}{cSec}{تمرين \theexo}}{\end{BOX}}
\newcommand{\ABC}[3]{\\[2pt]\begin{tabular}{p{0.28\textwidth}p{0.28\textwidth}p{0.28\textwidth}}(جـ) \ #3 & (ب) \ #2 & (أ) \ #1\end{tabular}}
\pretocmd{\section}{\Needspace{24\baselineskip}}{}{}
\newtcbox{\bxtbox}{on line,colback=cDef!12,colframe=cDef!12,boxrule=0pt,arc=1pt,boxsep=0pt,left=2pt,right=2pt,top=1pt,bottom=1pt}
\newcommand{\bxt}[1]{\foreignlanguage{english}{\bxtbox{\ar{#1}}}}
\SetHeader{\small\bfseries\color{cSec}تمارين البكالوريا: النهايات والاستمرارية والاشتقاقية}{\small\color{cSec}السنة الثالثة ثانوي -- الشعب العلمية}
\newenvironment{FIG}{\par\begin{center}\begin{LTR}}{\end{LTR}\end{center}}
\begin{document}
\begin{LTR}
\begin{tcolorbox}[enhanced,colback=cSec,colframe=cSec,arc=3mm,halign=center,boxrule=0pt,top=12pt,bottom=12pt,before upper={\selectlanguage{arabic}}]
{\color{white}\fontsize{21}{27}\selectfont\bfseries تمارين البكالوريا في النهايات والاستمرارية والاشتقاقية}\\[8pt]
{\color{white}\large -- بدون حلول (للتدرّب) -- الدوال كثيرات الحدود، الناطقة والصماء}\\[5pt]
{\color{white}\large علوم تجريبية -- تقني رياضي}
\end{tcolorbox}
\end{LTR}

\section{تمارين البكالوريا -- شعبة علوم تجريبية}
% ===================== بكالوريا 2014 -- الموضوع الثاني =====================
\begin{tamrin}{بكالوريا 2014 -- الموضوع الثاني -- علوم تجريبية}
\textbf{I.} لتكن $g$ الدالة العددية المعرفة على $\R$ كما يلي: $g(x)=2x^{3}-4x^{2}+7x-4$.
\QQ{1}
\begin{sq}
\qa احسب $\lim\limits_{x\to-\infty}g(x)$ و $\lim\limits_{x\to+\infty}g(x)$.
\qb ادرس اتجاه تغيّر الدالة $g$ على $\R$ ثم شكّل جدول تغيّراتها.
\end{sq}
\QQ{2}
\begin{sq}
\qa بيّن أن المعادلة $g(x)=0$ تقبل حلا وحيدا $\alpha$ حيث $0{,}7<\alpha<0{,}8$.
\qb استنتج حسب قيم العدد الحقيقي $x$ إشارة $g(x)$.
\end{sq}
\textbf{II.} نعتبر الدالة العددية $f$ المعرفة على $\R$ كما يلي: $f(x)=\dfrac{x^{3}-2x+1}{2x^{2}-2x+1}$، و $(C_f)$ تمثيلها البياني في المستوي المنسوب إلى معلم متعامد ومتجانس $(O;\vec i,\vec j)$.
\QQ{1} احسب $\lim\limits_{x\to-\infty}f(x)$ و $\lim\limits_{x\to+\infty}f(x)$.
\QQ{2}
\begin{sq}
\qa بيّن أنه من أجل كل $x$ من $\R$: $f(x)=\dfrac12(x+1)+\dfrac{1-3x}{2(2x^{2}-2x+1)}$.
\qb استنتج أن المنحنى $(C_f)$ يقبل مستقيما مقاربا مائلا $(\Delta)$ يطلب تعيين معادلة له.
\qc ادرس الوضع النسبي للمنحنى $(C_f)$ و $(\Delta)$.
\end{sq}
\QQ{3}
\begin{sq}
\qa بيّن أنه من أجل كل $x$ من $\R$: $f'(x)=\dfrac{x\,g(x)}{(2x^{2}-2x+1)^{2}}$ حيث $f'$ مشتقة الدالة $f$.
\qb استنتج إشارة $f'(x)$ حسب قيم $x$ ثم شكّل جدول تغيّرات الدالة $f$ (نأخذ $f(\alpha)\approx-0{,}1$).
\end{sq}
\QQ{4} احسب $f(1)$ ثم حل في $\R$ المعادلة $f(x)=0$.
\QQ{5} أنشئ المستقيم $(\Delta)$ والمنحنى $(C_f)$.
\QQ{6} لتكن $h$ الدالة العددية المعرفة على $\R$ بـ: $h(x)=\dfrac{x^{3}-4x^{2}+2x-1}{2x^{2}-2x+1}$، و $(C_h)$ تمثيلها البياني.
\begin{sq}
\qa تحقّق أنه من أجل كل $x$ من $\R$: $h(x)=f(x)-2$.
\qb استنتج أن $(C_h)$ هو صورة $(C_f)$ بتحويل نقطي بسيط يطلب تعيينه، ثم أنشئ $(C_h)$.
\end{sq}
\end{tamrin}

% ===================== بكالوريا 2009 -- الموضوع الأول =====================
\begin{tamrin}{بكالوريا 2009 -- الموضوع الأول -- علوم تجريبية}
\textbf{I.} $f$ دالة معرّفة على المجال $I=]-\infty;-1[\,\cup\,]-1;0]$ بـ: $f(x)=-x+\dfrac{4}{x+1}$، و $(C_f)$ تمثيلها البياني في مستوٍ منسوب إلى معلم متعامد ومتجانس $(O;\vec i,\vec j)$ كما هو مبيّن في الشكل الآتي.
\begin{GR}{-8}{1}{-5.5}{5}{0.6}
\Ticks{-7,-6,-5,-4,-3,-2,-1}{-5,-4,-3,-2,-1,1,2,3,4}
\draw[cThm,very thick,domain=-8:1] plot(\x,{-\x});
\draw[cThm,very thick,dashed](-1,-5.5)--(-1,5);
\draw[black,very thick,domain=-8:-1.15,samples=150] plot(\x,{-\x+4/(\x+1)});
\draw[black,very thick,domain=-0.6:0,samples=60] plot(\x,{-\x+4/(\x+1)});
\end{GR}
\QQ{1}
\begin{sq}
\qa احسب نهايات $f$ عند الحدود المفتوحة لـ $I$.
\qb بقراءة بيانية ودون دراسة اتجاه $f$، شكّل جدول تغيّراتها.
\end{sq}
\QQ{2} $g$ دالة معرّفة على $[0;+\infty[$ بـ: $g(x)=x+\dfrac{4}{x+1}$، و $(C_g)$ تمثيلها البياني.
\begin{sq}
\qa احسب نهاية $g$ عند $+\infty$.
\qb تحقّق من أن $(C_g)$ يقبل مستقيما مقاربا مائلا $(\Delta)$ عند $+\infty$ يطلب تعيين معادلة له.
\qc ادرس تغيّرات $g$.
\end{sq}
\textbf{II.} $k$ دالة معرّفة على $\R-\{-1\}$ بـ: $k(x)=|x|+\dfrac{4}{x+1}$.
\QQ{1}
\begin{sq}
\qa احسب $\lim\limits_{h\to0^{+}}\dfrac{k(h)-k(0)}{h}$ و $\lim\limits_{h\to0^{-}}\dfrac{k(h)-k(0)}{h}$.
\qb ماذا تستنتج؟ أعط تفسيرا هندسيا لهذه النتيجة.
\end{sq}
\QQ{2} اكتب معادلتي نصفي المماسين $(\Delta_1)$ و $(\Delta_2)$ عند النقطة التي فاصلتها $x_0=0$.
\QQ{3} ارسم $(\Delta_1)$ و $(\Delta_2)$ و $(C_k)$.
\end{tamrin}

% ===================== بكالوريا 2008 -- الموضوع الأول =====================
\begin{tamrin}{بكالوريا 2008 -- الموضوع الأول -- علوم تجريبية}
المنحنى $(C)$ الآتي هو التمثيل البياني للدالة العددية $g$ المعرّفة على $I=]-1;+\infty[$ بـ: $g(x)=x^{3}+3x^{2}+3x-1$.
\begin{GR}{-1.5}{2}{-2.5}{3.5}{1}
\Ticks{-1,1,2}{-2,-1,1,2,3}
\draw[blue,very thick,domain=-1:0.74,samples=100] plot(\x,{\x*\x*\x+3*\x*\x+3*\x-1});
\end{GR}
\QQ{1} بقراءة بيانية:
\begin{sq}
\qa شكّل جدول تغيّرات الدالة $g$ وحدّد $g(0)$ وإشارة $g(0{,}5)$.
\qb علّل وجود عدد حقيقي وحيد $\alpha\in]0;0{,}5[$ يحقق $g(\alpha)=0$.
\qc استنتج إشارة $g(x)$ على المجال $I$.
\end{sq}
\QQ{2} $f$ دالة معرّفة على $]-1;+\infty[$ بـ: $f(x)=\dfrac{x^{3}+3x^{2}+3x+2}{(x+1)^{2}}$، و $(\Gamma)$ تمثيلها البياني.
\begin{sq}
\qa تحقّق أنه من أجل كل $x\in I$: $f'(x)=\dfrac{g(x)}{(x+1)^{3}}$.
\qb عيّن دون حساب $\lim\limits_{x\to\alpha}\dfrac{f(x)-f(\alpha)}{x-\alpha}$ وفسّر النتيجة بيانيا.
\qc جد $\lim\limits_{x\to-1}f(x)$ و $\lim\limits_{x\to+\infty}\left[f(x)-(x+1)\right]$، فسّر النتيجتين بيانيا.
\qd شكّل جدول تغيّرات $f$.
\end{sq}
\QQ{3} نأخذ $\alpha=0{,}26$.
\begin{sq}
\qa عيّن مدوّر $f(\alpha)$ إلى $10^{-2}$.
\qb ارسم المنحنى $(\Gamma)$.
\end{sq}
\end{tamrin}

\section{تمارين البكالوريا -- شعبة تقني رياضي}
% ===================== بكالوريا 2017 -- الموضوع الأول =====================
\begin{tamrin}{بكالوريا 2017 -- الموضوع الأول -- تقني رياضي}
\textbf{I.} نعتبر الدالة العددية $g$ المعرّفة على $\R$ كما يلي: $g(x)=x^{3}+6x+12$.
\QQ{1} ادرس اتجاه تغيّر الدالة $g$.
\QQ{2} بيّن أن المعادلة $g(x)=0$ تقبل حلا وحيدا $\alpha$، $\alpha\in]-1{,}48;-1{,}47[$، ثم استنتج إشارة $g(x)$.

\textbf{II.} نعتبر الدالة العددية $f$ المعرّفة على $\R$ كما يلي: $f(x)=\dfrac{x^{3}-6}{x^{2}+2}$، و $(C_f)$ تمثيلها البياني في المستوي المنسوب إلى معلم متعامد ومتجانس $(O;\vec i,\vec j)$.
\QQ{1}
\begin{sq}
\qa احسب $\lim\limits_{x\to-\infty}f(x)$ و $\lim\limits_{x\to+\infty}f(x)$.
\qb بيّن أنه من أجل كل عدد حقيقي $x$: $f'(x)=\dfrac{x\,g(x)}{(x^{2}+2)^{2}}$، ثم ادرس اتجاه تغيّر الدالة $f$ وشكّل جدول تغيّراتها.
\end{sq}
\QQ{2}
\begin{sq}
\qa بيّن أن المستقيم $(\Delta)$ ذا المعادلة $y=x$ مقارب مائل للمنحنى $(C_f)$.
\qb ادرس وضعية المنحنى $(C_f)$ بالنسبة إلى المستقيم $(\Delta)$.
\end{sq}
\QQ{3} بيّن أن: $f(\alpha)=\dfrac32\alpha$، ثم استنتج حصرا للعدد $f(\alpha)$.
\QQ{4} ارسم المستقيم $(\Delta)$ والمنحنى $(C_f)$.
\end{tamrin}

% ===================== بكالوريا 2009 -- الموضوع الأول =====================
\begin{tamrin}{بكالوريا 2009 -- الموضوع الأول -- تقني رياضي}
$f$ دالة معرّفة على $]-1;+\infty[$ بـ: $f(x)=x-\dfrac{2}{\sqrt{x+1}}$، و $(C_f)$ منحنى $f$ في معلم متعامد ومتجانس $(O;\vec i,\vec j)$.
\QQ{1} ادرس تغيّرات الدالة $f$.
\QQ{2}
\begin{sq}
\qa بيّن أن $(C_f)$ يقبل مقاربين أحدهما $(D):\ y=x$.
\qb ادرس الوضعية النسبية للمنحنى $(C_f)$ و $(D)$.
\end{sq}
\QQ{3}
\begin{sq}
\qa بيّن أن المنحنى $(C_f)$ يقطع محور الفواصل في نقطة وحيدة فاصلتها $x_0$ حيث $1{,}3<x_0<1{,}4$.
\qb اكتب معادلة $(\Delta)$ مماس المنحنى $(C_f)$ في نقطة تقاطع $(C_f)$ مع حامل محور التراتيب.
\qc ارسم $(\Delta)$ و $(C_f)$ في نفس المعلم.
\end{sq}
\QQ{4} $g$ دالة معرّفة على المجال $]-1;+\infty[$ بـ: $g(x)=|f(x)|$، وليكن $(C_g)$ منحنى الدالة $g$ في نفس المعلم.
\begin{sq}
\qa بيّن كيف يمكن إنشاء $(C_g)$ انطلاقا من $(C_f)$ ثم ارسمه.
\qb ناقش بيانيا وحسب قيم الوسيط الحقيقي $m$ عدد وإشارة حلول المعادلة: $g(x)=m^{2}$.
\end{sq}
\end{tamrin}

\section{تمارين بمستوى البكالوريا}
% ===================== تمرين 1 =====================
\begin{tamrin}{تمرين بمستوى البكالوريا: دالة ناطقة، محور تناظر ومناقشة بيانية}
نعتبر $f$ الدالة المعرّفة على $\R-\{-1;3\}$ كما يلي: $f(x)=\dfrac{4x^{2}-8x}{x^{2}-2x-3}$، وليكن $(C_f)$ تمثيلها البياني في المستوي المنسوب إلى معلم متعامد ومتجانس $(O;\vec i,\vec j)$.
\QQ{1} عيّن الأعداد الحقيقية $a$ ، $b$ و $c$ بحيث يكون من أجل كل $x\neq-1$ و $x\neq3$: \ $f(x)=a+\dfrac{b}{x+1}+\dfrac{c}{x-3}$.
\QQ{2} احسب نهايات الدالة $f$ عند الأطراف المفتوحة لمجموعة التعريف.
\QQ{3} اكتب معادلة كل مستقيم من المستقيمات المقاربة لـ $(C_f)$.
\QQ{4} ادرس وضعية المنحنى $(C_f)$ بالنسبة إلى المستقيم المقارب الأفقي.
\QQ{5} ادرس اتجاه تغيّر الدالة $f$ ثم شكّل جدول تغيّراتها.
\QQ{6} عيّن $A$ و $B$ نقطتي تقاطع المنحنى $(C_f)$ مع حامل محور الفواصل.
\QQ{7} عيّن معادلتي المماسين $(\Delta)$ و $(\Delta')$ للمنحنى $(C_f)$ عند النقطتين $A$ و $B$، ثم عيّن إحداثيي نقطة تقاطع $(\Delta)$ و $(\Delta')$.
\QQ{8} أثبت أن المستقيم ذا المعادلة $x=1$ محور تناظر لـ $(C_f)$.
\QQ{9} احسب $f(-3)$ و $f(-2)$ ثم ارسم $(\Delta)$ ، $(\Delta')$ و $(C_f)$.
\QQ{10} ناقش بيانيا وحسب قيم الوسيط الحقيقي $m$ عدد وإشارة حلول كل من المعادلات:
\[f(x)=\sqrt{m}\qquad\qquad f(x)=\sqrt{(m+1)^{2}}\qquad\qquad (4+m)x^{2}-(8+2m)x-3m=0\]
\end{tamrin}

% ===================== تمرين 2 =====================
\begin{tamrin}{تمرين بمستوى البكالوريا: دالة بقيمة مطلقة، مقاربات مائلة ونقطة انعطاف}
نعتبر $f$ الدالة المعرّفة على $\R-\{-2;2\}$ كما يلي: $f(x)=|x+2|+\dfrac{4x}{x^{2}-4}$، وليكن $(C_f)$ تمثيلها البياني في المستوي المنسوب إلى معلم متعامد ومتجانس $(O;\vec i,\vec j)$.
\QQ{1} احسب $f(0)$ ، $f(-3)$ ، $f(-4)$.
\QQ{2} ادرس تغيّرات الدالة $f$.
\QQ{3}
\begin{sq}
\qa برهن على وجود عدد حقيقي $\alpha$ حيث $\alpha\in]-4;-3[$ و $f(\alpha)=0$.
\qb برهن أن المعادلة $f(x)=0$ تقبل حلا وحيدا $\beta\in]-2;2[$.
\end{sq}
\QQ{4} عيّن معادلات المستقيمات المقاربة العمودية والمائلة للمنحنى $(C_f)$.
\QQ{5}
\begin{sq}
\qa أوجد معادلة المماس $(\Delta)$ للمنحنى $(C_f)$ عند النقطة ذات الفاصلة $x_0=0$.
\qb ادرس وضعية المنحنى $(C_f)$ بالنسبة إلى المستقيم $(\Delta)$ على المجال $]-2;2[$ مفسّرا النتيجة هندسيا.
\end{sq}
\QQ{6} أنشئ المستقيمات المقاربة والمنحنى $(C_f)$.
\QQ{7} ناقش بيانيا وحسب قيم الوسيط الحقيقي $m$ عدد وإشارة حلول المعادلة: \ $|x+2|-2m=1-\dfrac{4x}{x^{2}-4}$.
\end{tamrin}

% ===================== تمرين 3 =====================
\begin{tamrin}{تمرين بمستوى البكالوريا: عائلة دوال $f_m$، مركز تناظر ومناقشة -- شعبة رياضيات}
\textbf{I.} نعتبر الدالة $f_m$ للمتغير الحقيقي $x$ والوسيط الحقيقي $m$ المعرّفة على $\R-\{3\}$ بـ: $f_m(x)=\dfrac{3x^{2}-mx}{x-3}$، وليكن $(C_m)$ المنحنى البياني الممثل للدالة $f_m$ في مستوٍ منسوب إلى معلم متعامد ومتجانس $(O;\vec i,\vec j)$.
\QQ{1} بيّن أنه توجد نقطة وحيدة ثابتة $A$ تنتمي إلى جميع المنحنيات $(C_m)$.
\QQ{2} ما هي قيم $m$ التي تجعل الدالة $f_m$ تتغيّر دوما في نفس الاتجاه؟
\QQ{3} ما هي قيم $m$ التي تجعل الدالة $f_m$ تقبل قيمة حدية محلية عظمى وقيمة حدية محلية صغرى؟

\textbf{II.} نعتبر الدالة $g$ المعرّفة بـ: $g(x)=f_5(x)$ (أي $m=5$)، و $(C)$ تمثيلها البياني في المعلم السابق.
\QQ{1} ادرس تغيّرات الدالة $g$.
\QQ{2} برهن أنه توجد نقطتان من المنحنى $(C)$ يكون عندهما ميل المماس $-2$، ثم اكتب معادلتي المماسين.
\QQ{3} عيّن الأعداد الحقيقية $a$ ، $b$ و $c$ بحيث من أجل كل $x\in D_g$: \ $g(x)=ax+b+\dfrac{c}{x-3}$.
\QQ{4} برهن أن المنحنى $(C)$ يقبل مستقيمين مقاربين ثم عيّن $\omega$ نقطة تقاطعهما.
\QQ{5} بيّن أن النقطة $\omega$ مركز تناظر للمنحنى $(C)$.
\QQ{6} أنشئ المنحنى $(C)$ ومستقيماته المقاربة في المعلم السابق.
\QQ{7} ناقش بيانيا وحسب قيم الوسيط الحقيقي $\lambda$ عدد وإشارة حلول المعادلة: \ $3x^{2}-(2\lambda+5)x+6\lambda=0$.

\textbf{III.} نعتبر الدالة $\varphi$ المعرّفة بـ: $\varphi(x)=\dfrac{3x^{2}-5|x|}{|x|-3}$.
\begin{itemize}[label=$\clubsuit$]
\item عيّن مجموعة تعريف الدالة $\varphi$.
\item ادرس قابلية اشتقاق الدالة $\varphi$ عند النقطة ذات الفاصلة $x_0=0$.
\item ادرس شفعية الدالة $\varphi$ (زوجية أم فردية).
\item برهن أن $\varphi(x)=g(x)$ على مجال $I$ يطلب تعيينه، ثم استنتج كيفية رسم منحنى الدالة $\varphi$.
\end{itemize}
\end{tamrin}

% ===================== تمرين 4 =====================
\begin{tamrin}{تمرين بمستوى البكالوريا: دالة مساعدة، حصر $f(\alpha)$ ومماسات موازية}
\textbf{I.} $g$ دالة معرّفة على $\R$ بـ: $g(x)=x^{3}+3x+8$.
\QQ{1} ادرس تغيّرات الدالة $g$ على $\R$ ثم شكّل جدول تغيّراتها.
\QQ{2} بيّن أن المعادلة $g(x)=0$ تقبل حلا وحيدا $\alpha$ في المجال $]-1{,}6;-1{,}5[$.
\QQ{3} حدّد حسب قيم $x$ إشارة $g(x)$.

\textbf{II.} $f$ دالة معرّفة بـ: $f(x)=\dfrac{x^{3}-4}{x^{2}+1}$، وليكن $(C_f)$ تمثيلها في معلم متعامد ومتجانس.
\QQ{1} عيّن $D_f$ مجموعة تعريف الدالة $f$.
\QQ{2} ادرس نهايات الدالة $f$ عند أطراف مجموعة تعريفها.
\QQ{3} بيّن أنه من أجل كل $x$ من $D_f$: \ $f'(x)=\dfrac{x\,g(x)}{(x^{2}+1)^{2}}$.
\QQ{4} ادرس اتجاه تغيّر الدالة $f$ ثم شكّل جدول تغيّراتها.
\QQ{5} برهن أن $f(\alpha)=\dfrac32\alpha$، ثم عيّن حصرا للعدد $f(\alpha)$.
\QQ{6} عيّن الأعداد الحقيقية $a$ ، $b$ و $c$ بحيث يكون $f(x)=ax+\dfrac{bx+c}{x^{2}+1}$ لكل $x$ من $\R$، ثم استنتج أن المنحنى $(C_f)$ له مستقيم مقارب مائل $(\Delta)$ يطلب تعيين معادلته.
\QQ{7} ادرس وضعية المنحنى $(C_f)$ بالنسبة إلى $(\Delta)$.
\QQ{8} عيّن نقاط تقاطع المنحنى $(C_f)$ مع حاملي المحورين.
\QQ{9} برهن أن المنحنى $(C_f)$ يقبل مماسين موازيين للمستقيم $(\Delta)$.
\QQ{10} احسب $f(-1)$ ، $f(1)$ ، $f(2)$ ثم ارسم $(\Delta)$ و $(C_f)$.

\textbf{III.} لتكن $h$ الدالة المعرّفة على $\R$ بـ: $h(x)=\left[f(x)\right]^{2}$.
\begin{itemize}[label=$\spadesuit$]
\item أعط عبارة $h'(x)$ بدلالة $f'(x)$ و $f(x)$ من أجل كل $x$ من $\R$.
\item ادرس إشارة $h'(x)$ من أجل كل $x$ من $\R$ ثم استنتج اتجاه تغيّر الدالة $h$ مشكّلا جدول تغيّراتها.
\end{itemize}
\end{tamrin}

% ===================== تمرين 5 =====================
\begin{tamrin}{تمرين بمستوى البكالوريا: مقارب مائل، مناقشة $|f(x)|$ ونقطة زاوية}
\textbf{I.} لتكن $g$ الدالة المعرّفة على $\R$ بـ: $g(x)=x^{3}+ax^{2}+bx+c$، و $(C_g)$ تمثيلها البياني.
\begin{itemize}[label=$\spadesuit$]
\item عيّن الأعداد الحقيقية $a$ ، $b$ ، $c$ حيث: $(C_g)$ يشمل النقطتين $A(1;0)$ ، $B(0;-5)$ ويقبل مماسا عند النقطة ذات الفاصلة $x_0=-1$ معادلته $y=-2x-6$.
\end{itemize}
نفرض فيما يلي أن: $a=3$ ، $b=1$ ، $c=-5$.
\QQ{1} ادرس تغيّرات الدالة $g$ على $\R$ ثم شكّل جدول تغيّراتها.
\QQ{2} احسب $g(1)$ ثم استنتج حسب قيم $x$ إشارة $g(x)$.

\textbf{II.} $f$ الدالة المعرّفة على $\R-\{-1\}$ بـ: $f(x)=\dfrac{x^{3}+3x^{2}+5x+5}{(x+1)^{2}}$ ، $(C_f)$ تمثيلها البياني في المستوي المنسوب إلى معلم متعامد ومتجانس $(O;\vec i,\vec j)$.
\QQ{1} احسب نهايات الدالة $f$ عند الأطراف المفتوحة لمجموعة التعريف.
\QQ{2} جد الأعداد الحقيقية $\alpha$ ، $\beta$ ، $\gamma$ ، $\lambda$ حتى يكون من أجل كل عدد حقيقي $x$ من $D_f$: \[f(x)=\alpha x+\beta+\frac{\gamma x+\lambda}{(x+1)^{2}}\]
\QQ{3} استنتج وجود مستقيم مقارب مائل $(\Delta)$ للمنحنى $(C_f)$ يطلب تعيين معادلته، ثم ادرس الوضع النسبي بينهما.
\QQ{4} بيّن أنه من أجل كل $x\neq-1$: \ $f'(x)=\dfrac{g(x)}{(x+1)^{3}}$.
\QQ{5} عيّن دون حساب: $\displaystyle\lim_{x\to1}\left(\frac{f(x)-f(1)}{x-1}\right)$ مفسّرا النتيجة بيانيا.
\QQ{6} ادرس اتجاه تغيّر الدالة $f$ وشكّل جدول تغيّراتها.
\QQ{7} بيّن أن المنحنى $(C_f)$ يقبل مماسا $(T)$ موازيا للمستقيم $(\Delta)$ يطلب كتابة معادلته. برّر أن المعادلة $f(x)=0$ تقبل حلا وحيدا $\theta$ في المجال $]-2;-1[$.
\QQ{8} ادرس نقاط تقاطع المنحنى $(C_f)$ مع حاملي المحورين ثم أنشئ بدقة $(\Delta)$ ، $(T)$ و $(C_f)$.
\QQ{9} ناقش بيانيا وحسب قيم الوسيط الحقيقي $\sigma$ عدد وإشارة حلول المعادلة: \ $|f(x)|=-\dfrac12\sigma+1$.

\textbf{III.} لتكن $h$ الدالة المعرّفة على $\R$ بـ: $h(x)=\dfrac{|x|^{3}+3x^{2}+5|x|+5}{\left(|x|+1\right)^{2}}$ ، $(C_h)$ تمثيلها البياني.
\begin{itemize}[label=$\star$]
\item بيّن أن $h$ دالة زوجية.
\item ادرس استمرارية الدالة $h$ عند النقطة ذات الفاصلة $0$.
\item بيّن أن: $\displaystyle\lim_{\substack{x\to0\\x<0}}\frac{h(x)-h(0)}{x}\neq\lim_{\substack{x\to0\\x>0}}\frac{h(x)-h(0)}{x}$ ثم فسّر النتيجة بيانيا.
\item اشرح كيفية رسم المنحنى $(C_h)$ انطلاقا من $(C_f)$ بعد نزع القيمة المطلقة لـ $h$.
\end{itemize}
\end{tamrin}

% ===================== تمرين 6 =====================
\begin{tamrin}{تمرين بمستوى البكالوريا: دالة صماء، مقارب مائل ومناقشة $g(x)=m^{2}$}
نعتبر $f$ الدالة المعرّفة على المجال $]-1;+\infty[$ بـ: $f(x)=x-\dfrac{2}{\sqrt{x+1}}$، و $(C)$ تمثيلها البياني.
\QQ{1} ادرس تغيّرات الدالة $f$ على المجال $]-1;+\infty[$.
\QQ{2} بيّن أن $(C)$ يقبل مستقيمين مقاربين أحدهما المستقيم $(\Delta)$ ذو المعادلة $y=x$، ثم ادرس الوضع النسبي بين $(C)$ و $(\Delta)$.
\QQ{3} بيّن أن المنحنى $(C)$ يقطع حامل محور الفواصل في نقطة وحيدة فاصلتها $x_0$ حيث: $1{,}3<x_0<1{,}4$.
\QQ{4} عيّن معادلة المماس $(T)$ للمنحنى $(C)$ عند نقطة تقاطعه مع حامل محور التراتيب.
\QQ{5} ارسم $(T)$ ، $(\Delta)$ و $(C)$.
\QQ{6} $g$ الدالة المعرّفة على $]-1;+\infty[$ بـ: $g(x)=|f(x)|$، و $(C_g)$ منحناها البياني.
\begin{itemize}[label=$\star$]
\item بيّن كيفية إنشاء $(C_g)$ انطلاقا من $(C)$ ثم أنشئه في المعلم السابق.
\item ناقش بيانيا حسب قيم الوسيط الحقيقي $m$ عدد وإشارة حلول المعادلة: $g(x)=m^{2}$.
\end{itemize}
\end{tamrin}

% ===================== تمرين 7 =====================
\begin{tamrin}{تمرين بمستوى البكالوريا: دالة ناطقة، دالة زوجية وعائلة منحنيات -- شعبة رياضيات}
\textbf{I.} نعتبر $f$ الدالة المعرّفة بـ: $f(x)=\dfrac{x^{2}+ax+b}{x^{2}+x-6}$.
\QQ{1} حدّد مجموعة تعريف الدالة $f$.
\QQ{2} جد العددين الحقيقيين $a$ و $b$ بحيث المنحنى $(C_f)$ يشمل النقطة $A(0;1)$ ويقبل مماسا معامل توجيهه $-\dfrac74$ عند النقطة ذات الفاصلة $x_0=1$.

\textbf{II.} لتكن الدالة $g$ المعرّفة على $\R-\{-3;2\}$ بـ: $g(x)=\dfrac{x^{2}+5x-6}{x^{2}+x-6}$.
\QQ{1} احسب نهايات الدالة $g$ عند حدود مجالات مجموعة التعريف.
\QQ{2} ادرس اتجاه تغيّر الدالة $g$ على مجموعة تعريفها.
\QQ{3} اكتب معادلة المماس $(T)$ عند النقطة ذات الفاصلة $x_0=1$.
\QQ{4} جد معادلات المستقيمات المقاربة للمنحنى $(C_g)$.
\QQ{5} ادرس وضعية المنحنى $(C_g)$ بالنسبة إلى المستقيم $(\Delta)$ ذي المعادلة: $y=1$.
\QQ{6} ارسم المنحنى $(C_g)$ والمستقيمات المقاربة.
\QQ{7} ناقش بيانيا حسب قيم الوسيط الحقيقي $m$ عدد وإشارة حلول المعادلة: \ $(m+1)x^{2}+(m+5)x-6m-6=0$.

\textbf{III.} لتكن الدالة $h$ المعرّفة بـ: $h(x)=\dfrac{x^{2}+5|x|-6}{x^{2}+|x|-6}$.
\QQ{1} جد مجموعة تعريف الدالة $h$.
\QQ{2} برهن أن الدالة $h$ زوجية وفسّر النتيجة بيانيا.
\QQ{3} اكتب $h(x)$ دون رمز القيمة المطلقة.
\QQ{4} عبّر عن $h$ بدلالة $g$ على مجال يطلب تعيينه.
\QQ{5} بيّن كيف يمكن إنشاء $(C_h)$ انطلاقا من $(C_g)$.

\textbf{IV.} لتكن الدالة $k_m$ المعرّفة بـ: $k_m(x)=\dfrac{x^{2}+5mx-6}{x^{2}+mx-6}$ ، حيث $m$ وسيط حقيقي.
\begin{itemize}[label=$\star$]
\item بيّن أن جميع المنحنيات $(C_{k_m})$ تشمل ثلاث نقاط ثابتة يطلب تعيين إحداثياتها.
\item ما هو المنحنى $(C_{k_m})$ الذي يشمل النقطة $B(-2;3)$؟
\end{itemize}
\end{tamrin}

% ===================== تمرين 8 =====================
\begin{tamrin}{تمرين بمستوى البكالوريا: قراءة بيانية، ثلاثة مقاربات ومناقشة بمستقيمات موازية}
\textbf{I.} المنحنى $(C_g)$ الآتي هو التمثيل البياني للدالة $g$ المعرّفة على $\R$ كما يلي: $g(x)=ax^{4}+bx^{2}+cx+d$ حيث $a$ ، $b$ ، $c$ ، $d$ أعداد حقيقية ثابتة (المستقيم المرسوم هو مماس $(C_g)$ عند النقطة ذات الفاصلة $0$).
\begin{GR}{-2.5}{2.5}{-2.5}{5.5}{0.75}
\Ticks{-2,-1,1,2}{-2,-1,1,2,3,4,5}
\draw[cThm,thick,domain=-2.5:2.5] plot(\x,{2*\x+1});
\draw[cEx,very thick,domain=-2.5:2.5,samples=200] plot(\x,{\x^4-2*\x*\x+2*\x+1});
\node[cEx] at (-1.6,3.6){$(C_g)$};
\end{GR}
\QQ{1} بقراءة بيانية حدّد $g(2)$ ، $g(1)$ ، $g'(0)$ ، $g(0)$.
\QQ{2} اعتمادا على ما سبق تحقق أن: $g(x)=x^{4}-2x^{2}+2x+1$.
\QQ{3} بيّن أن المعادلة $g(x)=0$ تقبل حلين مختلفين $\alpha$ و $\beta$ حيث: $-1{,}69<\alpha<-1{,}68$ و $-0{,}38<\beta<-0{,}37$.
\QQ{4} استنتج إشارة $g(x)$ على $\R$.

\textbf{II.} $f$ دالة معرّفة على $\R-\{-1;1\}$ بـ: $f(x)=x-\dfrac{1}{x^{2}-1}$، وليكن $(C_f)$ تمثيلها في معلم متعامد ومتجانس.
\QQ{1} احسب نهايات الدالة $f$ عند حدود مجموعة تعريفها.
\QQ{2} برهن أن المنحنى $(C_f)$ يقبل $3$ مستقيمات مقاربة يطلب تعيين معادلاتها.
\QQ{3} ادرس وضعية المنحنى $(C_f)$ بالنسبة إلى المقارب المائل $(\Delta)$.
\QQ{4} برهن أن الدالة $f$ قابلة للاشتقاق من أجل كل $x$ من $D_f$ وأن: $f'(x)=\dfrac{g(x)}{(x^{2}-1)^{2}}$، ثم استنتج اتجاه تغيّر $f$.
\QQ{5} أعط حصرا لـ $f(\alpha)$ و $f(\beta)$.
\QQ{6} تحقق أن المنحنى $(C_f)$ يقطع حامل محور الفواصل في نقطة فاصلتها $x_0$ حيث: $1{,}32<x_0<1{,}33$.
\QQ{7} بيّن أن المنحنى $(C_f)$ يقبل مماسا وحيدا $(T)$ موازيا للمستقيم $(\Delta)$ يطلب تعيين معادلته.
\QQ{8} أنشئ $(\Delta)$ ، $(T)$ و $(C_f)$.
\QQ{9} ناقش بيانيا وحسب قيم الوسيط $m$ عدد وإشارة حلول المعادلة: $mx^{2}+1-m=0$، ثم تحقق من ذلك جبريا.

\textbf{III.} لتكن $h$ الدالة المعرّفة على $\R-\{-1;1\}$ بـ: $h(x)=f(-x)$.
\begin{itemize}[label=$\star$]
\item باستعمال مشتقة مركّب دالتين ادرس اتجاه تغيّر الدالة $h$.
\end{itemize}
\end{tamrin}

% ===================== تمرين 9 =====================
\begin{tamrin}{تمرين بمستوى البكالوريا: قراءة بيانية، منحنى مقارب (قطع مكافئ) وانسحاب}
\textbf{I.} $g$ دالة معرّفة على $\R$ كما يلي: $g(x)=x^{3}+ax^{2}+bx+c$ حيث $a$ ، $b$ ، $c$ أعداد حقيقية ثابتة، والمنحنى $(C_g)$ الآتي هو التمثيل البياني لها:
\begin{GR}{-3}{2.5}{-5.5}{1.5}{0.8}
\Ticks{-2,-1,1,2}{-5,-4,-3,-2,-1,1}
\draw[cThm,thick,domain=-3:2.5] plot(\x,{-\x-4});
\draw[cDef,very thick,domain=-3:2.5,samples=200] plot(\x,{\x*\x*\x+\x*\x-\x-4});
\node[cDef] at (1.9,0.8){$(C_g)$};
\end{GR}
\QQ{1} عيّن قيم $a$ ، $b$ ، $c$ إذا علمت أن المنحنى $(C_g)$ يقبل مماسا معادلته $y=-x-4$ عند النقطة ذات الفاصلة $x_0$ ويقبل مماسا أفقيا عند النقطة التي فاصلتها $-1$.
\QQ{2} بقراءة بيانية شكّل جدول تغيّرات الدالة $g$ وحدّد $g(0)$ وإشارة $g(2)$.
\begin{itemize}[label=\textcircled{\scriptsize S}]
\item علّل وجود عدد حقيقي $\alpha$ من المجال $]0;2[$ يحقق $g(\alpha)=0$، ثم تحقق أن: $\alpha\in]1{,}4;1{,}5[$.
\item عيّن إشارة $g(x)$ على $\R$.
\end{itemize}

\textbf{II.} $f$ دالة معرّفة على $\R-\{-1\}$ بـ: $f(x)=\dfrac{x^{3}-x^{2}-4x+4}{x+1}$، و $(C_f)$ تمثيلها البياني في المستوي المنسوب إلى معلم متعامد ومتجانس $(O;\vec i,\vec j)$.
\QQ{1} احسب نهايات الدالة $f$ عند الأطراف المفتوحة لمجالي تعريفها.
\QQ{2} بيّن أنه من أجل كل عدد حقيقي $x$ من $\R-\{-1\}$: \ $f'(x)=\dfrac{2g(x)}{(x+1)^{2}}$.
\QQ{3} عيّن $\displaystyle\lim_{x\to\alpha}\frac{f(x)-f(\alpha)}{x-\alpha}$ دون حسابها مفسّرا النتيجة بيانيا.
\QQ{4} استنتج إشارة $f'(x)$ حسب قيم $x$ ثم شكّل جدول تغيّرات الدالة $f$.

ليكن $(\Gamma)$ التمثيل البياني للدالة $k$ المعرّفة على $\R$ بـ: $k(x)=x^{2}-2x-2$.
\QQ{5} احسب $\displaystyle\lim_{x\to\pm\infty}\left[f(x)-k(x)\right]$ وفسّر النتيجة المحصّل عليها بيانيا، ثم ادرس الوضع النسبي بين $(C_f)$ و $(\Gamma)$.
\QQ{6} تحقق أن: $f(\alpha)=\dfrac{9}{\alpha+1}-2\alpha-1$، ثم جد حصرا للعدد $f(\alpha)$.
\QQ{7} احسب $f(1)$ ثم حل في $\R$ المعادلة $f(x)=0$ مستنتجا نقاط تقاطع المنحنى $(C_f)$ مع حامل محور الفواصل.
\QQ{8} أنشئ $(C_f)$ و $(\Gamma)$ في نفس المعلم.

\textbf{III.}
\QQ{1} لتكن $h$ الدالة المعرّفة على $\R-\{-2\}$ بـ: $h(x)=\dfrac{x^{3}+2x^{2}-5x-4}{x+2}$، و $(C_h)$ تمثيلها البياني.
\begin{itemize}[label=\textcircled{\scriptsize S}]
\item تحقق أنه من أجل كل $x$ من $\R-\{-2\}$: $h(x)=f(x+1)-2$.
\item استنتج أن $(C_h)$ هو صورة $(C_f)$ بانسحاب يطلب تعيينه.
\end{itemize}
\QQ{2} لتكن $L$ الدالة المعرّفة بـ: $L(x)=-|f(x)|$، استنتج كيفية رسم $(C_L)$ انطلاقا من $(C_f)$.
\QQ{3} لتكن $M$ الدالة المعرّفة بـ: $M(x)=\sqrt{f(x)}$.
\begin{itemize}[label=\textcircled{\scriptsize S}]
\item جد مجموعة تعريف الدالة $M$.
\item جد عبارة $M'$ بدلالة $f'$ و $f$ (حيث $M'$ هي الدالة المشتقة للدالة $M$).
\end{itemize}
\end{tamrin}

% ===================== تمرين 10 =====================
\begin{tamrin}{تمرين بمستوى البكالوريا: من جدول التغيّرات إلى العبارة، مركز تناظر ومناقشة $f(x)=|\lambda|$}
لتكن $f$ دالة قابلة للاشتقاق على كل مجال من مجموعة تعريفها، لها جدول التغيّرات التالي:
\begin{TAB}
\tkzTabInit[lgt=2,espcl=2.6]{$x$/1,$f'(x)$/1,$f(x)$/2}{$-\infty$,$\frac12$,$1$,$\frac32$,$+\infty$}
\tkzTabLine{,+,z,-,d,-,z,+,}
\tkzTabVar{-/$-\infty$,+/$1$,-D+/$-\infty$/$+\infty$,-/$3$,+/$+\infty$}
\end{TAB}
تُكتب عبارة $f$ على الشكل: $f(x)=ax+b+\dfrac{c}{x-1}$ حيث $a$ ، $b$ ، $c$ أعداد حقيقية.
\QQ{1} احسب $f'(x)$.
\QQ{2} اعتمادا على جدول تغيّرات الدالة $f$:
\begin{itemize}[label=\textcircled{\scriptsize S}]
\item عيّن الأعداد الحقيقية $a$ ، $b$ ، $c$.
\item عيّن $\lim\limits_{\substack{x\to1\\x<1}}f(x)$ و $\lim\limits_{\substack{x\to1\\x>1}}f(x)$ وفسّر النتيجة بيانيا.
\item قارن بين صورتي العددين $\dfrac12$ و $\dfrac34$ بالدالة $f$ معلّلا إجابتك.
\end{itemize}
\QQ{3} نأخذ فيما يلي: $a=1$ ، $b=1$ ، $c=\dfrac14$، وليكن $(C_f)$ تمثيلها البياني في المستوي المنسوب إلى معلم متعامد ومتجانس $(O;\vec i,\vec j)$.
\begin{itemize}[label=\textcircled{\scriptsize S}]
\item بيّن أن المنحنى $(C_f)$ له مستقيم مقارب مائل $(\Delta)$ يطلب تعيين معادلته.
\item ادرس وضعية المنحنى $(C_f)$ بالنسبة إلى المستقيم $(\Delta)$.
\item أثبت أن النقطة $\omega(1;2)$ مركز تناظر للمنحنى $(C_f)$.
\item عيّن نقطة تقاطع المنحنى $(C_f)$ مع حامل محور الفواصل.
\end{itemize}
\QQ{4} أنشئ $(C_f)$ و $(\Delta)$.
\QQ{5} ناقش بيانيا وحسب قيم الوسيط الحقيقي $\lambda$ عدد حلول المعادلة: $f(x)=|\lambda|$.
\end{tamrin}

% ===================== تمرين 11 =====================
\begin{tamrin}{تمرين بمستوى البكالوريا: دالة صماء، نقطة انعطاف ومناقشة بمستقيم دوّار -- شعبة رياضيات}
\textbf{I.} نعتبر الدالة العددية $g$ المعرّفة على $\R$ بـ: $g(x)=\dfrac12x^{2}-\sqrt{x^{2}+1}+1$.
\QQ{1} تحقق أنه من أجل كل $x$ من $\R$: \ $g'(x)=\left(\sqrt{x^{2}+1}-1\right)\dfrac{x}{\sqrt{x^{2}+1}}$.
\QQ{2} بيّن أنه من أجل كل $x$ من $\R$: $\sqrt{x^{2}+1}-1\geqslant0$، ثم شكّل جدول تغيّرات الدالة $g$.
\QQ{3} استنتج أنه من أجل كل عدد حقيقي $x$: $g(x)\geqslant0$.

\textbf{II.} نعتبر الدالة العددية $f$ المعرّفة على $\R$ بـ: \ $\left\{\begin{array}{ll}f(x)=\dfrac{1-\sqrt{x^{2}+1}}{x}, & x\neq0\\[3mm] f(0)=0\end{array}\right.$ \ و $(C_f)$ تمثيلها البياني في المستوي المنسوب إلى معلم متعامد ومتجانس $(O;\vec i,\vec j)$.
\QQ{1} أثبت أن الدالة $f$ مستمرة عند $0$.
\QQ{2}
\begin{itemize}[label=\textcircled{\scriptsize S}]
\item احسب $\displaystyle\lim_{x\to0}\frac{f(x)}{x}$ ثم استنتج أن الدالة $f$ تقبل الاشتقاق عند $0$.
\item أثبت أن المستقيم $(T)$ ذا المعادلة $y=-\dfrac12x$ مماس للمنحنى عند النقطة ذات الفاصلة $x_0=0$.
\item أثبت أنه من أجل كل عدد حقيقي $x$ غير معدوم: $f(x)+\dfrac12x=\dfrac{g(x)}{x}$.
\item ادرس الوضعية النسبية بين المنحنى $(C_f)$ والمستقيم $(T)$، ثم استنتج أن $(C_f)$ يقبل نقطة انعطاف يطلب تعيين إحداثييها.
\end{itemize}
\QQ{3} احسب نهاية الدالة $f$ عند $+\infty$ و $-\infty$ وفسّر النتائج بيانيا.
\QQ{4} بيّن أنه من أجل كل عدد حقيقي $x$ غير معدوم: \ $f'(x)=-\left(\sqrt{x^{2}+1}-1\right)\dfrac{1}{x^{2}\sqrt{x^{2}+1}}$.
\QQ{5} استنتج اتجاه تغيّر الدالة $f$ على $\R$ ثم شكّل جدول تغيّراتها.
\QQ{6} أنشئ $(C_f)$ ، $(T)$ والمستقيمات المقاربة.
\QQ{7} $m$ عدد حقيقي من المجال $\left[-\dfrac12;0\right]$، ناقش بيانيا وحسب قيم الوسيط الحقيقي $m$ عدد وإشارة حلول المعادلة: $f(x)=mx$.
\end{tamrin}

% ===================== تمرين 12 =====================
\begin{tamrin}{تمرين بمستوى البكالوريا: دالة مساعدة، مقارب مائل ومناقشة $f(x)=m$ و $-|f|$}
\textbf{I.} $g$ دالة معرّفة على $\R$ بـ: $g(x)=-2x^{3}+4x^{2}-7x+4$.
\QQ{1} ادرس تغيّرات الدالة $g$ على $\R$ ثم شكّل جدول تغيّراتها.
\QQ{2} بيّن أن المعادلة $g(x)=0$ تقبل حلا وحيدا $\alpha$ في المجال $]0{,}78;0{,}79[$.
\QQ{3} استنتج حسب قيم $x$ إشارة $g(x)$.

\textbf{II.} $f$ دالة معرّفة على $\R$ بـ: $f(x)=\dfrac{3x^{3}+2x^{2}-8x+4}{2x^{2}-2x+1}$، وليكن $(C_f)$ تمثيلها في المستوي المنسوب إلى المعلم المتعامد والمتجانس $(O;\vec i,\vec j)$.
\QQ{1} احسب نهايات الدالة $f$ عند أطراف مجموعة تعريفها.
\QQ{2} بيّن أنه من أجل كل $x$ من $D_f$: \ $f'(x)=\dfrac{-3x\,g(x)}{(2x^{2}-2x+1)^{2}}$.
\QQ{3} ادرس اتجاه تغيّر الدالة $f$ ثم شكّل جدول تغيّراتها.
\QQ{4} بيّن أنه من أجل كل عدد حقيقي $x$: \ $f(x)=\dfrac12(3x+5)+\dfrac{3-9x}{2\left(2x^{2}-2x+1\right)}$.
\QQ{5} استنتج أن المنحنى $(C_f)$ يقبل مستقيما مقاربا مائلا $(\Delta)$ يطلب تعيين معادلته.
\QQ{6} ادرس وضعية المنحنى $(C_f)$ بالنسبة إلى $(\Delta)$.
\QQ{7} بيّن أن المعادلة $f(x)=0$ تقبل حلا وحيدا $x_0$ حيث: $-2{,}18<x_0<-2{,}17$.
\QQ{8} أنشئ $(\Delta)$ و $(C_f)$ (نأخذ $f(\alpha)\approx0{,}6$).
\QQ{9} ناقش بيانيا وحسب قيم الوسيط الحقيقي $m$ عدد وإشارة حلول المعادلة: \[3x^{3}+2(1-m)x^{2}+2(m-4)x+4-m=0\]

\textbf{III.} لتكن $h$ الدالة المعرّفة على $\R$ بـ: $h(x)=-|f(x)|$.
\begin{itemize}[label=$\spadesuit$]
\item اكتب $h(x)$ دون رمز القيمة المطلقة.
\item استنتج كيفية إنشاء $(C_h)$ انطلاقا من $(C_f)$، ثم أنشئه في المعلم السابق.
\end{itemize}
\end{tamrin}

% ===================== تمرين 13 =====================
\begin{tamrin}{تمرين بمستوى البكالوريا: قراءة بيانية، دالة ناطقة ودراسة $f\left(\frac1x\right)$}
لتكن $g$ الدالة المعرّفة على $\R$ بـ: $g(x)=x^{3}+ax+b$، وليكن $(C_g)$ تمثيلها البياني و $(T)$ المماس لـ $(C_g)$ عند النقطة ذات الفاصلة $0$.
\begin{GR}{-1.5}{2.5}{-4.5}{2.8}{0.85}
\Ticks{-1,1,2}{-4,-3,-2,-1,1,2}
\draw[cThm,thick,domain=-0.85:1.6] plot(\x,{3*\x-2});
\draw[cDef,very thick,domain=-0.75:1.25,samples=150] plot(\x,{\x*\x*\x+3*\x-2});
\node[cDef] at (0.5,2.0){$(C_g)$}; \node[cThm] at (1.95,1.9){$(T)$};
\end{GR}
بقراءة بيانية:
\QQ{1} عيّن $g(0)$ ، $g'(0)$ و $\displaystyle\lim_{x\to0}\left(\frac{g'(x)-3}{x}\right)$.
\QQ{2} علّل وجود عدد حقيقي $\alpha$ من المجال $]0{,}5;0{,}7[$ يحقق $g(\alpha)=0$.
\QQ{3} عيّن حسب قيم العدد الحقيقي $x$ إشارة $g(x)$ ثم استنتج إشارة $g(-x)$.
\QQ{4} عيّن العددين الحقيقيين $a$ و $b$.

فيما يلي نضع: $a=3$ و $b=-2$. \ $f$ الدالة المعرّفة على $\R$ بـ: $f(x)=\dfrac{1-x^{3}}{1+x^{2}}$.
\QQ{1} احسب نهايات الدالة $f$ عند أطراف مجموعة تعريفها.
\QQ{2}
\begin{itemize}[label=$\spadesuit$]
\item بيّن أنه من أجل كل عدد حقيقي $x$: $f'(x)=\dfrac{x\,g(-x)}{(1+x^{2})^{2}}$.
\item حدّد اتجاه تغيّر الدالة $g$ ثم شكّل جدول تغيّراتها.
\end{itemize}
\QQ{3} عيّن دون حساب النهاية الآتية مفسّرا النتيجة بيانيا: $\displaystyle\lim_{x\to-\alpha}\left(\frac{f(x)-f(-\alpha)}{x+\alpha}\right)$.
\QQ{4}
\begin{itemize}[label=$\clubsuit$]
\item بيّن أن المستقيم $(\Delta)$ ذا المعادلة $y=-x$ مقارب مائل لـ $(C_f)$.
\item ادرس الوضع النسبي لـ $(C_f)$ و $(\Delta)$.
\end{itemize}
\QQ{5} بيّن أن: $f(-\alpha)=\dfrac32\alpha$، ثم استنتج حصرا للعدد $f(-\alpha)$.
\begin{itemize}[label=$\star$]
\item أنشئ $(C_f)$ و $(\Delta)$ (نأخذ $f(-\alpha)=0{,}7$).
\item ناقش بيانيا وحسب قيم الوسيط الحقيقي $m$ عدد وإشارة حلول المعادلة: \ $x^{2}(x-2m)=2m+1$.
\end{itemize}
\QQ{6} $h$ الدالة المعرّفة على $\R$ بـ: $h(x)=-f(|-x|)$.
\begin{itemize}[label=$\blacklozenge$]
\item بيّن أن $h$ دالة زوجية.
\item اشرح كيف يمكن إنشاء منحنى الدالة $h$ انطلاقا من $(C_f)$ ثم ارسمه.
\end{itemize}
\QQ{7} نعتبر $F$ الدالة المعرّفة على $\R^{*}$ بـ: $F(x)=f\left(\dfrac1x\right)$. دون تعيين عبارة $F$ ادرس تغيّرات الدالة $F$ وشكّل جدول تغيّراتها.
\end{tamrin}

% ===================== تمرين 14 =====================
\begin{tamrin}{تمرين بمستوى البكالوريا: من المنحنى إلى العبارة، نقطة انعطاف ونقطة زاوية}
$(C_f)$ التمثيل البياني للدالة $f$ المعرّفة على $D_f$ في المستوي المنسوب إلى المعلم المتعامد والمتجانس $(O;\vec i,\vec j)$. $(\Delta)$ ذو المعادلة $y=x$ المستقيم المقارب المائل لـ $(C_f)$ بجوار $+\infty$ و $-\infty$، و $x=2$ معادلة المستقيم المقارب الموازي لحامل محور التراتيب، $(T)$ المماس للمنحنى $(C_f)$ عند النقطة $B(3;1)$ والمار من النقطة $C(2;-3)$، و $A(1;1)$ نقطة من المنحنى $(C_f)$ كما هو موضّح في الشكل أدناه:
\begin{GR}{-3}{5}{-5}{5}{0.68}
\Ticks{-2,-1,1,3,4}{-4,-3,-2,-1,1,2,3,4}
\draw[cThm,dashed,thick](2,-5)--(2,5);
\draw[cEx,thick,domain=-3:5] plot(\x,{\x});
\draw[cRem,thick,domain=1.5:4] plot(\x,{4*\x-11});
\draw[cDef,very thick,domain=-3:1.62,samples=200] plot(\x,{\x+(1-\x)/((\x-2)*(\x-2))});
\draw[cDef,very thick,domain=2.4:5,samples=200] plot(\x,{\x+(1-\x)/((\x-2)*(\x-2))});
\fill (1,1) circle(2pt) (3,1) circle(2pt) (2,-3) circle(2pt);
\node[above] at (1,1.1){$A$}; \node[right] at (3.05,0.8){$B$}; \node[left] at (2,-3){$C$};
\node[cDef] at (-2.3,-1.9){$(C_f)$}; \node[cEx] at (4.4,3.6){$(\Delta)$}; \node[cRem] at (3.4,3.9){$(T)$};
\end{GR}
\textbf{I.} بقراءة بيانية:
\QQ{1} عيّن $D_f$ مجموعة تعريف الدالة $f$.
\QQ{2} عيّن: $\lim\limits_{x\to+\infty}f(x)$ ، $\lim\limits_{x\to-\infty}f(x)$ ، $\lim\limits_{x\to2}f(x)$ ، $\lim\limits_{x\to\pm\infty}\left(f(x)-x\right)$ ، $f(1)$ ، $f(3)$ ، $f'(1)$ ، $f'(3)$.
\QQ{3} عيّن معادلة المماس $(T)$ للمنحنى $(C_f)$ عند النقطة $B$.
\QQ{4} نفرض أنه من أجل كل $x$ من $D_f$: $f(x)=ax+\dfrac{bx+c}{(x-2)^{2}}$ حيث $a$ ، $b$ ، $c$ أعداد حقيقية.
\begin{itemize}[label=\textcircled{\scriptsize S}]
\item احسب $f'(x)$ عبارة المشتقة الأولى للدالة $f$ بدلالة $a$ ، $b$ و $c$.
\item باستعمال المعطيات السابقة عيّن الأعداد $a$ ، $b$ و $c$.
\end{itemize}

\textbf{II.}
\QQ{1} تحقق أنه من أجل كل $x\in\R-\{2\}$: \ $f(x)=\dfrac{x^{3}-4x^{2}+3x+1}{(x-2)^{2}}$.
\QQ{2} بيّن أنه من أجل كل عدد حقيقي $x$ من $\R-\{2\}$: \ $f'(x)=\dfrac{(x-1)(x^{2}-5x+8)}{(x-2)^{3}}$.
\QQ{3} استنتج إشارة $f'(x)$ حسب قيم $x$ ثم شكّل جدول تغيّرات الدالة $f$.
\QQ{4} استنتج وجود مستقيم مقارب مائل $(\Delta)$ للمنحنى $(C_f)$ يطلب تعيين معادلته، ثم ادرس الوضع النسبي بينهما.
\QQ{5} بيّن أن المنحنى $(C_f)$ يقبل مماسا $(T')$ موازيا للمستقيم $(\Delta)$ يطلب كتابة معادلته.
\QQ{6} بيّن أنه من أجل كل $x$ من $\R-\{2\}$: $f''(x)=\dfrac{-2x-2}{(x-2)^{4}}$، ثم ادرس إشارة $f''$؛ ماذا تستنتج بالنسبة إلى $(C_f)$؟
\QQ{7} ناقش بيانيا حسب قيم الوسيط $m$ عدد وإشارة حلول المعادلة: \ $x^{3}-(m+4)x^{2}+(3+4m)x+(1-4m)=0$.

\textbf{III.} لتكن $h$ الدالة المعرّفة على $\R-\{2\}$ بـ: \ $h(x)=\left\{\begin{array}{ll}x-3+\dfrac{1}{x-2}, & x<3\\[3mm] f(x), & x\geqslant3\end{array}\right.$ \ و $(C_h)$ تمثيلها البياني.
\begin{itemize}[label=\textcircled{\scriptsize S}]
\item ادرس استمرارية الدالة $h$ عند $x_0=3$. هل تقبل الدالة $h$ الاشتقاق عند $x_0=3$؟
\item عيّن معادلات المستقيمات المقاربة لـ $(C_h)$.
\end{itemize}
\end{tamrin}

\end{document}
